\section{Problem and method}\label{sec:method}
\paragraph{Setting.}
A leaderboard is a score matrix $X\in[0,100]^{M\times n}$ of $M$ models on $n$ tasks. A rule $f$ maps each score vector to a value (score rules) or each matrix to an outcome (rank rules), and ranks models by it. Fix a target model $i$ with scores $x=X_{i\cdot}$. An \emph{intervention} is a vector $y$ with $x_t\le y_t\le c_t$, where $c_t$ is a cap on task $t$. All other models stay fixed: no rival responds and no spillover moves other tasks. The target is \emph{first} if it is strictly ahead of every rival under $f$ by a margin $\varepsilon=10^{-3}$ points (for rank rules, strictly better outcome on a discrete level model in which a tie with a rival counts half).

\paragraph{Investment.}
The cost of $y$ is measured in leaderboard points, in three ways:
$K=|\{t:y_t>x_t\}|$ (tasks touched), $L_1=\sum_t(y_t-x_t)$ (total points) and $L_\infty=\max_t(y_t-x_t)$ (largest single-task gain; the uniform ``general improvement'' benchmark).
Each is the exact minimum of its own problem, so the minimiser of one is generally not that of another and they are never plotted as one strategy.
\emph{This is not a cost in data, compute or money.} The map from a data effort to points on a task is outside the model; the cap is where we encode what is plausible.

\paragraph{Caps.}
The cap is the main modelling assumption, so it is a swept parameter: \textsf{ceil} (100, a theoretical upper bound only), \textsf{front0}/\textsf{front2} (the best rival's score on that task plus 0 or 2 points), \textsf{gain5}/\textsf{gain10}/\textsf{gain20} (at most that many points per task), and \textsf{emp95}/\textsf{emp99} (per-task 95th/99th percentiles of fine-tune-minus-base differences observed on OLLB; 540 pairs). The empirical caps reflect general fine-tuning gains on six tasks, not the effect of task-specific data; on HELM only the other caps are used. An infeasible target (no raise within the cap reaches first place) is recorded as such and never dropped.

\paragraph{Rules.}
Score rules: arithmetic mean (the OLLB rule), geometric and harmonic mean (on scores shifted by 1, checked with other offsets), median, minimum, trimmed mean (drop the best and worst task; OLLB only), and weighted means as a sensitivity check. Rank rules: Borda, plurality (number of tasks on which a model is best), and the Condorcet requirement (the target must beat every rival on a majority of tasks). The last stands for Condorcet-consistent rules (Copeland, Schulze and others), not for Borda or Kemeny, and Kemeny is out of scope. Score rules differ in which part of the score vector matters (all of it, the middle, the minimum); rank rules ignore magnitude.

\paragraph{Exact solvers.}
For the mean the minimum has a closed form: with deficit $R=nb-\sum_t x_t$ and $b$ the best rival's value plus $\varepsilon$, $L_1=R$ whenever the headroom suffices, $K$ is the least number of the largest headrooms that cover $R$, and $L_\infty$ is a water level. The minimum and median also have closed forms for $K$, $L_1$ and $L_\infty$ (odd $n$; for even $n$ only $K$), and none of them is ``$n$ tasks'' or ``$\lceil n/2\rceil+1$ tasks'' in general, because tasks already above the winning level count (Appendix~\ref{app:theory}). Geometric and power means are solved by water-filling with bisection, the trimmed mean and rank rules by mixed-integer programmes solved with HiGHS \citep{huangfu2018highs} with lazily added rival constraints. Every solver is tested against brute force on small random instances, and the closed forms against solver output on every target at the ceiling cap: \num{theory.mismatch} mismatches in \num{theory.cells} cells (largest difference \num{theory.maxerr}); a numerical agreement, not a proof.

\paragraph{Targets and denominators.}
The targets are every model except the leader under the plain arithmetic mean: \num{helm.ceil.mean.ntar} of 79 on HELM and \num{ollb.ceil.mean.ntar} of 4496 on OLLB for score rules. Rank rules on OLLB are costly, so they are run on a rank-stratified sample of \num{ollb_sample.ceil.mean.ntar} targets, a deterministic grid (ranks 2--21 plus log-spaced ranks), not a population sample; it is labelled wherever it appears. A target already first under a non-mean rule has zero cost under it and stays in the comparison. Costs are summarised as quantiles over targets and, for capped runs, over feasible targets only, together with the feasible share (Wilson 95\% interval).

\paragraph{Uncertainty.}
We keep three kinds of interval apart. \emph{Wilson intervals} describe a share of targets. \emph{Bootstrap over targets} (2000 resamples, percentile) describes the sampling variability of a median over the target models; it is not evaluation noise. \emph{Draw quantiles} describe the spread over replicates for one target: evaluation noise (binomial standard errors per cell from the item counts), bootstrap over tasks, random task subsets, and random weights; these are the quantity of interest and not confidence intervals.
Evaluation-noise draws give a distribution of the re-optimised minimum cost. For ``how likely is a fixed investment to put the target first'' we replay a strategy chosen on the observed leaderboard (a uniform raise of $d$ points on every task up to the cap, or the same $L_1$ on the tasks with most headroom) against noise draws: each cell gets an independent Gaussian error with the binomial standard error of that cell (item counts, observed score), clipped to $[0,100]$; the success curves depend on this assumption. Wilson bands on replay curves are pointwise Monte Carlo intervals over the draws, and target-bootstrap bars resample the observed target models, not evaluation noise and not a population of models.
